\centerline{\bf \Huge Разнобой 9 класс}

\begin{flushright}
        \scriptsize{Но мы всё равно не умеем считать!\\ Правая голова Огр мага
}
\end{flushright}
 
\problem{1} 
Два приведённых квадратных трёхчлена $f(x)$ и $g(x)$ таковы, что каждый из них имеет по два корня, и выполняются равенства $f(1) = g(2)$ и $g(1) = f(2)$. Найдите сумму всех четырёх корней этих трёхченов.

\problem{2}
Петя и Миша стартуют по круговой дорожке из одной точки в направлении против часовой стрелки. Оба бегут с постоянными скоростями, скорость Миши на 2\% больше скорости Пети. Петя всё время бежит против часовой стрелки, а Миша может менять направление бега в любой момент, непосредственно перед которым он пробежал полуокружность или больше в одном направлении. Покажите, что пока Петя бежит первый круг, Миша может трижды, не считая момента старта, поравняться (встретиться или догнать) с ним.

\problem{3}
Вася записал в клетки таблицы $9 \times 9$ натуральные числа от $1$ до $81$ (в каждой клетке стоит по числу, все числа различны). Оказалось, что любые два числа, отличающихся на $3$, стоят в соседних по стороне клетках. Верно ли, что обязательно найдутся две угловых клетки, разность чисел в которых делится на $6$?

\problem{4}
На доске девять раз (друг под другом) написали некоторое натуральное число $N$. Петя к каждому из $9$ чисел приписал слева или справа одну ненулевую цифру; при этом все приписанные цифры различны. Какое наибольшее количество простых чисел могло оказаться среди $9$ полученных чисел?

